\chapter{Preliminaries}
\section{General Definitions}
This section introduces essential notation and mathematical concepts. \mycomment{2.1: Würde die einzelnen Subsections mit je einem Satz einleiten und ggf. ref-fen}

\subsection{Sets}
A \textit{set} is a collection of elements, denoted by capital letters (e.g., $A$, $S$, $X$). Elements within sets are represented by lower-case letters.

The following special sets are used throughout this thesis:
\begin{itemize}
    \item The empty set: $\emptyset$,
    \item The set of natural numbers: $\mathbb{N} \defeq \{0, 1, 2, \dots\}$,
    \item The set of natural numbers up to $n$: $\underline{n}\defeq\set{i\in\N\mid 0\leq i\leq n}$
    \item The set of integers: $\mathbb{Z} \defeq \{\dots, -2, -1, 0, 1, 2, \dots\}$,
    \item The set of real numbers: $\mathbb{R}$,
    \item Intervals: $[a, b] \defeq \{x \in \mathbb{R} \mid a \leq x \leq b\}$ and $(a, b) \defeq \{x \in \mathbb{R} \mid a < x < b\}$,
    \item The power set: $2^A$, which is the set of all subsets of $A$ \mycomment{das Komma weg?}
\end{itemize}

\subsection{Functions}
A \textit{function} is a mapping from each element of a set $X$ (called the \textit{domain}) to a single element in a set $Y$ (called the \textit{codomain}). Functions are typically denoted by lower-case letters, such as $f$. 

The function $f$ from $X$ to $Y$ is denoted as:
\[
f: X \to Y
\]

For an element $x \in X$, the function’s value at $x$ is denoted by $f(x) \in Y$. A function can also be defined explicitly, for example:
\[
f(x) \defeq x^2 + 1
\]
\mycomment{Warum ändert sich das Seiten layout alle 2 Seiten!??? Die Breite zum Seitenrand?}
A \textit{partial function} has the special codomain element $\bot$. If $f(x) = \bot$ for some $x$, the function is considered undefined at $x$.

In some cases, standard naming conventions for sets and functions are adjusted for clarity, especially when commonly accepted alternative symbols exist.

\subsection{Discrete Probability Distributions}
The following definitions on probability are based on Cramer \cite[Chapter 2]{cramer_statistik_2007}. Let $\Omega$ be a finite or countably\mycomment{besser: "countable"?} infinite set, referred to as the \textit{sample space}. Define a function $p: \Omega \to [0, 1]$ such that:
\[
\sum_{\omega \in \Omega} p(\omega) = 1
\]

A \textit{discrete probability distribution} is a function $\Pr: 2^\Omega \to [0, 1]$ defined as:
\[
\Pr(A) \defeq \sum_{\omega \in A} p(\omega) \quad \text{for} \quad A \in 2^\Omega
\]

The function $p$ is called the \textit{probability mass function} of $\Pr$. A probability distribution $\Pr$ is called \textit{deterministic} if $p(\omega)=1$ for some $\omega\in\Omega$.

The set $\mathcal{D}(\Omega)$ represents all discrete probability distributions over the sample space $\Omega$. The \textit{dirac distribution} for an element $x\in\Omega$ is defined by \[
\delta_x(A)\defeq\begin{cases}
    1 & \text{if } x \in A\\
    0 & \text{otherwise}\\
\end{cases}\]

\subsection{Probability spaces}
Not all probabilistic events that we are interested in are subsets of countable sets. For example, later in this thesis, we want to assign probabilities to subsets of the set of all infinite paths in a graph, which is an uncountable set. Therefore, we introduce \textit{probability spaces}.
\mycomment{Warum erst hier Def-enviornments genutzt und nicht schon bei Sets?}
\begin{mydefinition}{Probability space}{def:prob}
A probability space is a tuple $(\Omega, \mathcal{F}, \mu)$, where:
\begin{itemize}
    \item $\Omega \neq \emptyset$ is a (possibly uncountable) sample space,
    \item $\mathcal{F} \subseteq 2^\Omega$ is a \textit{event space}, \mycomment{event space genutzt bevor es definiert wurde (oder soll das hier die def sein? - dann lieber "is called a" schreiben)}
    \item $\mu: \mathcal{F} \to [0, 1]$ is a probability measure,
\end{itemize}
and it holds that:
\begin{enumerate}
    \item $\Omega \in \mathcal{F}$,
    \item $\mathcal{F}$ is closed under complementation: $\forall A \in \mathcal{F},\quad A^c = \Omega \setminus A \in \mathcal{F}$, \mycomment{für mich uneindeutig was hier $\in \mathcal{F}$ ist}
    \item $\mathcal{F}$ is closed under countable union: $\forall A_0, A_1, \dots \in \mathcal{F}, \quad \bigcup_{i \in \mathbb{N}} A_i \in \mathcal{F}$,
    \item $\mu(\Omega) = 1$,
    \item $\mu$ is \textit{countably additive}: For any disjoint sets $A_0, A_1, \dots \in \mathcal{F}$, where $A_i \cap A_j = \emptyset$ for $i \neq j$, it holds that:
    \[
    \mu\left(\bigcup_{i \in \mathbb{N}} A_i\right) = \sum_{i \in \mathbb{N}} \mu(A_i).
    \]
\end{enumerate}
An event space $\mathcal{F}$ is called a \textit{$\sigma$-algebra} if conditions 1–3 are satisfied, and a probability function $\mu$ on a $\sigma$-algebra $\mathcal{F}$ is called a \textit{probability measure} if conditions 4 and 5 are satisfied.

Let $\mathcal{C} \subseteq 2^\Omega$ be a collection of sets. The least $\sigma$-algebra containing $\mathcal{C}$ is called the \textit{$\sigma$-algebra generated by $\mathcal{C}$}, denoted by $\sigma(\mathcal{C})$. It exists, is unique, and corresponds to the closure of $\mathcal{C}$ under conditions 1–3. \mycomment{Quelle?}
\end{mydefinition}



\subsection{Languages}
The following definitions are taken from Staiger \cite{staiger_omega-languages_1997}. Let $\Sigma$ be a finite set of symbols, called the \textit{alphabet}. A \textit{word} over $\Sigma$ is a finite sequence of elements $a_0a_1\ldots$, where $a_i \in \Sigma$ for all $i \geq 0$. The empty word is denoted by $\varepsilon$. Sets of words are referred to as \textit{languages}. Given languages $L$ and $L'$ over $\Sigma$, the following operators are defined:
\begin{itemize}
    \item Concatenation: $LL' \defeq \{ ww' \mid w \in L, w' \in L' \}$
    \item $n$-fold concatenation: $L^n \defeq \underbrace{L L \ldots L}_{n \text{ times}} = \{ w_0w_1\ldots w_n \mid w_i \in L, i \in \underline{n} \},\quad n\in\N$
    \item Kleene star: $L^\ast \defeq \bigcup_{i \in \mathbb{N}} L^i$
    \item Infinite concatenation: $L^\omega \defeq \{ w_0w_1\ldots \mid w_i \in L, i \in \mathbb{N} \}$ (where $L \neq \emptyset$ or $\{\varepsilon\}$). Concatenation of $L^\omega$ with $L'$ on the right ($L^\omega L'$) is undefined.
    \item Finite or infinite concatenation: $L^{\leq\omega} \defeq L^\ast \cup L^\omega$
\end{itemize}

Correspondingly, the set of all finite words over an alphabet $\Sigma$ is $\Sigma^\ast$, the set of all infinite words is $\Sigma^\omega$, and the set of all finite and infinite words is $\Sigma^{\leq\omega}$.\mycomment{$\omega$ nur implizit definiert?}




\section{Stochastic Multi-Player Games}
This thesis uses the definition by Chen et al. \cite[Definition 1]{chen_automatic_2013}. In a stochastic multi-player game (SMG), the system is represented as a finite graph. Each vertex corresponds to a player and has an associated set of available actions. The game proceeds as follows: In each round, the active player (i.e., the player associated with the current state) selects an action at the current vertex. Subsequently, the next state is chosen probabilistically based on a distribution determined by the current state and the chosen action.
\vspace{2mm}

\begin{mydefinition}{Stochastic Multi-Player Game (SMG)}{def:smg}
A SMG $\mathcal{G}$ is a tuple $(\Pi, S, A, (S_i)_{i\in\Pi}, \Delta, AP, \chi)$, where:
\begin{itemize}
    \item $\Pi$ is a finite set of players,
    \item $S$ is a finite, non-empty set of states,
    \item $A$ is a finite, non-empty set of actions,
    \item $(S_i)_{i\in\Pi}$ is a partition of $S$,
    \item $\Delta: S \times A \to \mathcal{D}(S) \cup \{\bot\}$ is a partial transition function,
    \item $AP$ is a finite set of atomic propositions,
    \item $\chi: S \to 2^{AP}$ is a labelling function.
\end{itemize}
\end{mydefinition}

For each state $s \in S$, let $A(s) \defeq \{a \in A \mid \Delta(s, a) \neq \bot\}$ denote the set of available actions in state $s$. The current player, meaning the player $i \in \Pi$ such that $s \in S_i$, selects an action $a \in A(s)$, and the next state is drawn according to $\Delta(s, a)$.

A finite or infinite sequence of states and actions $s_0a_0s_1a_1...\in (SA)^\ast S\cup(SA)^\omega$ is called a path. It is called \textit{valid} if $a_i\in A(s_i)$ and $\Delta(s_i, a_i)(s_{i+1})>0$ . Finite paths are usually denoted by $\lambda=s_0a_0s_1a_1...$, infinite paths by $\omega$. The sequence of states in a path $\lambda$ is denoted by $st_\lambda\defeq s_0s_1...$, the set of all infinite paths is $\Omega_\mathcal{G}$ and the set of infinite paths starting in a state $s'\in S$ is $\Omega_{\mathcal{G},s'}$.

\begin{figure}[t]
    \centering
\begin{tikzpicture}[>=stealth, node distance=3cm, every node/.style={font=\small}]

% Nodes
\node[draw, minimum size=1cm, shape=rectangle] (S1) {$s_1$};
\node[draw, minimum size=1cm, regular polygon,regular polygon sides=3,right= 5cm of S1] (S2) {$s_2$};
\node[draw, minimum size=1cm, shape=circle, below= 3cm of $(S1)!0.5!(S2)$] (S3) {$s_3$};

\draw[-] (S1) -- node[below left, pos=0.5] {$a$} ($(S1)!0.5!(S3)$);
\filldraw[black] ($(S1)!0.5!(S3)$) circle (2pt); % Dot at split point
\draw[->] ($(S1)!0.5!(S3)$) -- node[below left, pos=0.5] {$\frac{2}{3}$} (S3);
\draw[->] ($(S1)!0.5!(S3)$) -> node[below, pos=0.75] {$\frac{1}{3}$} (S2);
\draw[-] (S2) -- node[below right, pos=0.5] {$a$} ($(S2)!0.5!(S3)$);
\filldraw[black] ($(S2)!0.5!(S3)$) circle (2pt); % Dot at split point
\draw[->] ($(S2)!0.5!(S3)$) -- node[below right, pos=0.5] {$\frac{1}{2}$} (S3);
\draw[->] ($(S2)!0.5!(S3)$) -- node[below, pos=0.75] {$\frac{1}{2}$} (S1);

\draw[->] (S1) to[bend left=30] node[above, pos=0.25] {$b$} 
node[above, pos=0.75] {$1$} coordinate[pos=0.5] (S1S2) (S2);
\filldraw[black] (S1S2) circle (2pt); % Dot halfway along the curve

\draw[->] (S2) -> node[above, pos=0.25] {$b$} node[above, pos=0.75] {$1$} coordinate[pos=0.5] (S2S1) (S1);
\filldraw[black] (S2S1) circle (2pt); % Dot halfway along the curve

\draw[->] (S3) edge[loop below]
coordinate[pos=0.5] (S3S3)
node[right, pos=0.2] {$c$}
node[left, pos=0.8] {$1$} (S3);
\filldraw[black] (S3S3) circle (2pt); % Dot halfway along the curve

\end{tikzpicture}
    \caption{Example of a SMG with three players and three actions}
    \label{fig:example_game_01}
\end{figure}


\begin{myexample}{Simple SMG}{ex:simple_smg}
Figure\~\ref{fig:example_game_01} shows an example of a SMG. It has three players: \textit{square}, \textit{triangle} and \textit{circle}. The outgoing edges correspond to the available action, e.g.,$A(s_1)=\set{a, b}$. The multiple ends of each edge represent the states that have a non-zero probability according to the transition function $\Delta$. They are labeled with their probabilities. E.g.,$\Delta(s_1, a)(s_2)=\frac{1}{3}$ and $\Delta(s_1, a)(s_3)=\frac{2}{3}$.
\end{myexample}

\subsection{Strategies}\label{def:strat}

To reason about nondeterministic behavior of multiple players, it is useful to introduce \textit{strategies}. A strategy is a function that describes the decision a player makes in each step.
\begin{mydefinition}{Strategy}{def:strat}
Let $\mathcal{G}=(\Pi, S, A, (S_i)_{i\in\Pi}, \Delta, AP, \chi)$ be a SMG. A strategy for a player $i\in\Pi$ is a function $\sigma_i: (SA)^\ast S_i\rightarrow\mathcal{D}(A)$, 
where $\sigma_i(\lambda s)$ is a probability distribution over $A(s)$ and $\lambda s\in (SA)^\ast S_i$.
\end{mydefinition}
A strategy $\sigma_i$ is called \textit{memoryless} if it is path-independent, this means $\sigma_i(\lambda s) = \sigma_i(\lambda' s)$ for all $\lambda$ and $\lambda'\in (SA)^\ast, s\in S_i$ and is called \textit{deterministic} if $\sigma_i(\lambda s)$ is deterministic.

$\Sigma_i$ denotes the set of all strategies for player $i\in\Pi$. A \textit{strategy profile} $\sigma=\sigma_0,...,\sigma_{|\Pi|-1}$ contains a strategy for each player.
\mytodo{example of strategy on \ref{fig:example_game_01}}


\subsection{Outcomes of Strategies}\label{sec:prop_measure}
An SMG $\mathcal{G}$ paired with a strategy profile becomes a purely probabilistic process, as all non-deterministic choices of players are determined. To formalize this, we introduce a $\sigma$-algebra over the paths in $\mathcal{G}$, following the definitions of Chatterjee \cite[Section 2.2.2]{Chatterjee07_StochasticOmegaRegularGames}.\mycomment{REF der Quelle funktioniert hier nicht}

Given a finite path $\lambda = s_0 a_0 s_1 a_1 \dots s_n$, the \textit{cone} of $\lambda$ is the set 
\[
\Cone(\lambda) \defeq \set{ s'_0 a'_0 s'_1 \dots \in \Omega_\mathcal{G} \mid s_i = s'_i \text{ and } a_i = a'_i, \; i \in \underline{n}}
\]
of infinite paths that have $\lambda$ as their prefix. Let $\mathcal{U} \defeq \left\{ \text{Cone}(\lambda) \mid \lambda \text{ is a finite path in } \mathcal{G} \right\}$, and let $\mathcal{F} \defeq \sigma(\mathcal{U})$ be the $\sigma$-algebra generated by $\mathcal{U}$.

\begin{mydefinition}{Strategy-Induced Probability Measure}{def:strat_prop}
Let $\mathcal{G} = (\Pi, S, A, (S_i)_{i \in \Pi}, \Delta, AP, \chi)$ be an SMG, $\sigma$ a strategy profile for $\mathcal{G}$, $s \in S$, and let $\mathcal{U}$ and $\mathcal{F}$ be defined as above. We define the function $\mu_s^\sigma: \mathcal{U} \rightarrow [0, 1]$ as follows:

Let $x \in \mathcal{U}$. If $x = \text{Cone}(s')$ with $s' \in S$, then
\[
\mu_s^\sigma(x) \defeq
\begin{cases}
1 & \text{if } s = s', \\
0 & \text{otherwise.}
\end{cases}
\]
Otherwise $x = \Cone(s_0 a_0 \dots s_n a_n s_{n+1})$ for some $n \geq 0$, then
\[
\mu_s^\sigma(Cone(s_0 a_0 \dots s_n a_n s_{n+1})) \defeq \mu_s^\sigma(Cone(s_0 a_0 s_1 a_1 \dots s_n)) \cdot \sigma_i(a_n) \cdot \Delta(s_n, a_n)(s_{n+1}),
\]
where $i\in\Pi$ s.t. $s_n \in S_i$.

$\mu_s^\sigma$ is a measure, and there is a unique extension of $\mu_s^\sigma$ on $\mathcal{F}$ as a probability measure $\Pr_s^\sigma$ (according to the Carath\'eodory Extension Theorem \cite[Theorem 11.2]{billingsley_probability_1995}). Thus $(\Omega_\mathcal{G}, \mathcal{F}, \Pr_s^\sigma)$ is a probability space.

$\Pr_s^\sigma$ is the \textit{strategy-induced probability measure} of $\mathcal{G}$ with the strategy profile $\sigma$ and initial state $s$.
\end{mydefinition}

Note that $\mathcal{F}$ does not include all subsets of $\Omega_\mathcal{G}$, but only those that are finite Boolean combinations of cone sets. This means there are properties of paths to which we cannot assign a probability. However, the properties we are interested in are all expressible as subsets of $\mathcal{F}$.


\section{Reachability}

A useful question about SMGs is whether a player or a set of players (a \textit{coalition}) can ensure that a set of target states is eventually reached. Here, we differentiate between cases:
\begin{enumerate}
    \item A player can ensure the target set is reached with a non-zero probability, we call this \textit{reachability}, or
    \item a player can ensure the target set is reached with probability $1$, we call this \textit{almost-sure reachability} (AS-reachability).
\end{enumerate}
In this section we will formalize these kinds of reachability.

Instead of reasoning about a coalition of players, we can reduce SMGs to two-player games. All states of the players in the coalition belong to player $0$, whereas all adversarial states belong to player $1$. It is easy to see that states are (almost-sure) reachable for a coalition of players if and only if they are (almost-sure) reachable for player $0$ in the reduced SMG, as the strategies of the players in the coalition can be combined into a strategy for player $0$.

Thus, throughout this section, let $\mathcal{G} = (\set{0, 1}, S, A, \set{S_0, S_1}, \Delta, AP, \chi)$ be a two-player SMG and $T \subseteq S$ be a set of target states. We define the set of infinite paths $\text{Reach}(T)$ that eventually lead to a state in $T$:
\[
    \Reach(T) \defeq \set{s_0 a_0 s_1 \dots \in \Omega_\mathcal{G} \mid \exists k \geq 0 \text{ s.t. } s_k \in T}
\]
Let $\mathcal{F}$ be the $\sigma$-algebra of $\mathcal{G}$ as defined in Section \ref{sec:prop_measure}. It holds that
\[
\Reach(T) = (SA)^\ast T (AS)^\omega = \bigcup_{w \in (SA)^\ast T} w (AS)^\omega = \bigcup_{w \in (SA)^\ast T} \Cone(w) \in \mathcal{F}
\]
as it is a countable union of cone sets. This means $\Reach(T)$ is measurable with $\Pr_s^{(\sigma_0, \sigma_1)}$ for $s \in S$ and any strategy profile $(\sigma_0, \sigma_1)$. This allows as to define reachability in terms of probabilities on $\mathcal{F}$:

\begin{mydefinition}{(Almost-Sure) Reachability}{def:reach}
Let $\mathcal{G} = (\set{0, 1}, S, A, \set{S_0, S_1}, \Delta, AP, \chi)$ be a two-player SMG and $T \subseteq S$ be a set of target states. $T$ is \textit{reachable} for player $0$ from an initial state $s$, if there exists a strategy $\sigma_0$ for player $0$ such that for all strategies $\sigma_1$ of player $1$, it holds that:
\[
\Pr_s^{(\sigma_0, \sigma_1)}(\Reach(T)) > 0
\]
$T$ is \textit{almost surely reachable} if it holds that:
\[
\Pr_s^{(\sigma_0, \sigma_1)}(\Reach(T)) = 1
\]
We denote the set of states from which $T$ is (almost surly) reachable for the player $i\in\set{0, 1}$ with:\begin{align*}
    \Reach_i(T, \mathcal{G})&\defeq\set{s\in S\mid\exists \sigma_i\forall\sigma_{1-i}: \Pr_s^{(\sigma_0, \sigma_1)}(\Reach(T)) > 0}\\
    \ASReach_i(T, \mathcal{G})&\defeq\set{s\in S\mid\exists \sigma_i\forall\sigma_{1-i}: \Pr_s^{(\sigma_0, \sigma_1)}(\Reach(T)) = 1}
\end{align*}
\end{mydefinition}

A useful result for reachability in SMGs is \textit{determinacy}, which states that either a player has a strategy to achieve their objective, or the opponent has a strategy to prevent them from doing so. We formally state this property in a way that will be useful later in the thesis; it can also be expressed more generally.
\begin{mytheorem}{Reachability on SMGs is determined}{theo:reach_det}\label{theo:reach_det}
Either player $i$ has a memoryless strategy $\sigma_i$, such that for all strategies $\sigma_{1-i}$ it holds that $\Pr_s^{(\sigma_0, \sigma_1)}(\Reach(T))=0$ or there is a $\delta > 0$ and a strategy $\sigma_{1-i}$ such that for all strategies $\sigma_i$, it holds that $\Pr_s^{(\sigma_0, \sigma_1)}(\Reach(T)) \geq \delta$.
\end{mytheorem}
This formulation corresponds to the result of Chatterjee et al. \cite{Chatterjee07_StochasticOmegaRegularGames}[Theorem 2], which establishes determinacy for the more general case of parity objectives.
